% Generated draft: requires mathematical and source review.
\documentclass[11pt]{article}
\usepackage[margin=25mm]{geometry}
\usepackage{fontspec}
\setmainfont{Latin Modern Roman}
\usepackage{amsmath,amssymb,mathtools,unicode-math}
\setmathfont{Latin Modern Math}
\usepackage{enumitem,needspace,xurl,hyperref}
\setlength{\parindent}{0pt}
\setlength{\parskip}{5pt}
\setlength{\emergencystretch}{3em}
\newcommand{\N}{\mathbb{N}}
\newcommand{\Z}{\mathbb{Z}}
\newcommand{\Q}{\mathbb{Q}}
\newcommand{\R}{\mathbb{R}}
\newcommand{\C}{\mathbb{C}}
\newcommand{\F}{\mathbb{F}}
\newcommand{\A}{\mathbb{A}}
\newcommand{\PP}{\mathbb{P}}
\newcommand{\E}{\mathbb{E}}
\newcommand{\eps}{\varepsilon}
\newcommand{\dd}{\,\mathrm{d}}
\begin{document}
\section*{UnsolvedMath ID: 30001457}
\textbf{Source problem code:} catalogue entry
\subsection{Definitions and mathematical statement}
Let \(A=\{a_1,\ldots,a_k\}\), and let \(A^{\Z}\) denote the set of bi-infinite words \(w=(w_n)_{n\in\Z}\) with letters in \(A\). The source asks for a characterization of cutting sequences of geodesics on translation surfaces obtained from regular \(2k\)-gons, together with a reconstruction of a geodesic from its sequence.

The following makes one conventional interpretation precise, conditionally on choices not specified in the record. For an integer \(k\geq 2\), let \(P\subset\R^2\) be a nondegenerate regular \(2k\)-gon. Identify each side with its opposite side by translation to obtain a translation surface \(S\), and label the \(k\) pairs of identified sides bijectively by \(a_1,\ldots,a_k\). Retain the polygonal presentation, so that crossings of the identified sides can be recorded. A translation surface is a surface with Euclidean coordinate charts away from finitely many distinguished points, whose transition maps are translations; a geodesic is locally straight in these charts.

For this conditional interpretation, let \(\mathcal G(S)\) consist of oriented unit-speed geodesics \(\gamma:\R\to S\) which avoid all images of polygon vertices, cross the identified sides transversely, and have their entire set of crossing times indexed as
\[
\cdots<t_{-1}<t_0=0<t_1<\cdots,
\qquad \lim_{n\to\pm\infty}t_n=\pm\infty.
\]
Thus a crossing is distinguished as index zero. Define the cutting sequence map
\[
C_S:\mathcal G(S)\longrightarrow A^{\Z},\qquad
C_S(\gamma)_n=a_i
\]
when the crossing at time \(t_n\) is through the identified side pair labelled \(a_i\).

\textbf{Conditional mathematical formulation.} For every such \(k\) and labelled surface \(S\), give necessary and sufficient conditions on \(w\in A^{\Z}\) for
\[
w\in C_S(\mathcal G(S)).
\]
For every realizable \(w\), determine the fiber
\[
C_S^{-1}(w)=\{\gamma\in\mathcal G(S):C_S(\gamma)=w\},
\]
and determine in what sense the geodesic can be recovered from \(w\), including whether it is unique under the intended equivalence of geodesics.

\textbf{Unresolved qualifications.} The source does not explicitly specify opposite-side identification, the permitted values of \(k\), the labelling convention, or the treatment of vertex encounters, singular geodesics, and trajectories along sides. The class \(\mathcal G(S)\) above is therefore a conditional regular-trajectory formulation, not a verified exhaustive interpretation of the source. The meaning of recovery is also unspecified: it could concern direction, the geometric trajectory, or a parametrized geodesic, possibly modulo time translation or surface symmetries. If the distinguished crossing is forgotten, coding is naturally considered modulo the index shift \((\sigma w)_n=w_{n+1}\). These choices require mathematical review before a definitive formulation can be accepted.
\subsection{Short English statement}
Characterize exactly which bi-infinite words encode geodesics on translation surfaces obtained from regular \(2k\)-gons, and determine what geometric information the word recovers. The displayed formulation conditionally uses opposite-side gluing and geodesics avoiding polygon vertices with transverse crossings indexed by \(\Z\); the source leaves these conventions and the meaning of unique recovery unresolved.
\subsection{Sources}
\begin{itemize}
\item Supplied local UnsolvedMath registry record, permanent ID \(30001457\), title \emph{Cutting Sequences on Regular Polygon Translation Surfaces}, source problem code \texttt{catalogue entry}, in the collection \emph{Oberwolfach Reports Open Problems}. The record gives proposed year 2010 and names no individual proposer.
\item Supplied report citation: \emph{Arbeitsgemeinschaft: Mathematical Billards} (2010), Oberwolfach Reports, DOI 10.4171/owr/2010/17, pp.\ 955--1015. \url{https://doi.org/10.4171/owr/2010/17}. The supplied background describes these reports as records of mathematical workshops and their open-problem discussions; no problem-specific page is supplied.
\item Supplied literature source labelled ``Johnson, square-tiled surfaces'': \url{https://arxiv.org/abs/1605.05659}.
\item Supplied literature source labelled ``Davis--Pasquinelli--Ulcigrai'': \url{https://arxiv.org/abs/1509.03905}.
\item Literature assessment dated exactly \texttt{2026-08-21}: the supplied curation labels the problem open and states that complete symbolic characterization and geometric reconstruction are known for several special translation surfaces, while the general classification and reconstruction are not settled for arbitrary rational polygonal billiards. This is supplied dated curation, not an independently verified assessment. Its reference to arbitrary rational polygonal billiards is broader than the stated regular-polygon problem and does not itself establish the status of that narrower problem. The source expressly requests independent verification of time-sensitive status and cited claims.
\item The registry labels the statement \texttt{corrected\_verified}, explaining only that TeX notation was enclosed in mathematical delimiters without changing the mathematical content. This label is supplied metadata and does not establish mathematical verification of the formulation.
\end{itemize}
\end{document}
